\documentclass[10pt]{article}
\usepackage[T1]{fontenc}
\usepackage[utf8]{inputenc}
\usepackage[french]{babel}
\usepackage{amsmath,amssymb,mathrsfs,mathpazo}
\usepackage{geometry}
\geometry{a4paper,margin=2cm}
\usepackage{graphicx,array,tabularx,tabularray,colortbl,multirow,multicol,booktabs,enumitem,pifont,ulem,fancybox,diagbox,cancel}
\usepackage{pstricks,pstricks-add,pst-plot,pst-tree,pst-node,pst-text,pst-func,pst-math,pst-bspline,pst-3dplot,pst-eucl,pst-solides3d}
\let\pscalculate\undefined
\usepackage{pst-bezier}
\usepackage{tikz,pgfplots,tkz-tab}
\usetikzlibrary{arrows,arrows.meta,patterns,calc,positioning,trees,decorations.markings,shapes.geometric,shapes.misc,shapes.arrows,intersections,decorations.pathmorphing,decorations.pathreplacing,backgrounds,shadows,babel}
\pgfplotsset{compat=1.18}
\usepgfplotslibrary{fillbetween}
\usepackage[np]{numprint}
\usepackage{listings,makecell,pas-tableur,twoopt}
\let\barre\undefined
\usepackage[pstricks]{tablvar}
\definecolor{darkgreen}{RGB}{0,100,0}
\providecommand{\red}{\color{red}}
\providecommand{\blue}{\color{blue}}
\providecommand{\green}{\color{green}}
\providecommand{\black}{\color{black}}
\providecommand{\white}{\color{white}}
\providecommand{\gray}{\color{gray}}
\providecommand{\cyan}{\color{cyan}}
\providecommand{\magenta}{\color{magenta}}
\providecommand{\yellow}{\color{yellow}}
\providecommand{\violet}{\color{violet}}
\providecommand{\orange}{\color{orange}}
\providecommand{\euro}{\texteuro}
\providecommand{\EUR}{\texteuro}
\providecommand{\ieme}{\textsuperscript{e}}
\providecommand{\degres}{\ensuremath{^\circ}}
\providecommand{\degree}{\ensuremath{^\circ}}
\providecommand{\vv}[1]{\overrightarrow{#1}}
\providecommand{\Coord}[1]{#1}
\providecommand{\decosix}{}
\providecommand{\starredbullet}{\ensuremath{\bullet}}
\pagestyle{empty}
\setlength{\parindent}{0pt}
\usepackage{adjustbox}\DeclareUnicodeCharacter{2236}{:}\DeclareUnicodeCharacter{2190}{\ensuremath{\leftarrow}}
\providecommand{\R}{\mathbb{R}}
\providecommand{\N}{\mathbb{N}}
\providecommand{\Z}{\mathbb{Z}}
\providecommand{\Q}{\mathbb{Q}}
\providecommand{\D}{\mathbb{D}}
\providecommand{\C}{\mathbb{C}}
\providecommand{\e}{{\mathrm{e}}}
\providecommand{\vect}[1]{\overrightarrow{#1}}
\providecommand{\vectt}[1]{\overrightarrow{\mathrm{#1}}}
\providecommand{\vvt}[1]{\overrightarrow{\mathrm{#1}}}
\providecommand{\np}[2][]{{#2~#1}}
\providecommand{\vv}[1]{{\overrightarrow{#1}}}
\renewcommand{\d}{{\mathrm{d}}}
\providecommand{\degres}{{^\circ}}
\providecommand{\degre}{{^\circ}}
\providecommand{\blue}{\color{blue}}
\providecommand{\red}{\color{red}}
\providecommand{\enskip}{\quad}
\providecommand{\llbracket}{\lbrack\!\lbrack}
\providecommand{\rrbracket}{\rbrack\!\rbrack}
\providecommand{\euro}{{\text{€}}}
\providecommand{\no}{{\text{n°}}}
\providecommand{\interval}[2]{{[#1;#2]}}
\providecommand{\intervFO}[2]{{[#1;#2[}}
\providecommand{\dotfill}{\ldots}

\providecommand{\ds}{\displaystyle}
\providecommand{\textding}[1]{#1}
\providecommand{\Oij}{\(\mathrm{O};\vec{\imath},\vec{\jmath}\)}
\providecommand{\Oijk}{\(\mathrm{O};\vec{\imath},\vec{\jmath},\vec{k}\)}
\providecommand{\Ouv}{\(\mathrm{O};\vec{u},\vec{v}\)}
\providecommand{\intervFF}[2]{\ensuremath{[#1;#2]}}
\renewcommand{\euro}{\ifmmode\text{\texteuro}\else\texteuro\fi}
\DeclareUnicodeCharacter{00B0}{\ifmmode^{\circ}\else\textdegree\fi}
\DeclareUnicodeCharacter{20AC}{\ifmmode\text{\texteuro}\else\texteuro\fi}
% Énoncés : Éducation nationale. Transcription et corrigés : APMEP, Denis Vergès et contributeurs. Découpage et classement : CoopMaths, Sébastien Lozano, Éric Elter et contributeurs. Adaptation de la présentation : Maths Explorer.
% CC BY-SA : https://coopmaths.fr/www/about/
% https://coopmaths.fr/alea/static/bac/2026/tex/bac_2026_06_sujet1_polynesie_3.tex
\begin{document}
\noindent\textit{APMEP / CoopMaths. CC BY-SA. Présentation adaptée par Maths Explorer.}\par\medskip
On considère la fonction $f$ définie sur $[0~;~+\infty[$ par~:

\[f(x) = \ln(3x^2 + 1)\]

On admet que la fonction $f$ est croissante sur $[0~;+\infty[$.

\subsection*{Partie A}

On considère la suite $(u_n)$ définie par~:

\[u_0 = 2 \quad \text{et pour tout entier } n \geqslant 0,\quad u_{n+1} = \ln \left(3u_n^2 + 1\right).\]

\begin{enumerate}
\item Calculer $u_1$.
\item La suite $(u_n)$ est représentée ci-dessous. Émettre des conjectures sur le sens de variation et la convergence de la suite $(u_n)$.

\begin{center}























\psset{xunit=0.8cm,yunit=1.6cm,arrowsize=2pt 3}
\begin{pspicture*}(-0.7,-0.4)(13,4.5)
\psgrid[gridlabels=0pt,subgriddiv=1,gridwidth=0.15pt]
\psaxes[linewidth=1.25pt,Dx=2]{->}(0,0)(0,0)(13,4.5)
\psdots(0,2)(1,2.565)(2,3.032)(3,3.353)(4,3.547)(5,3.657)(6,3.717)
(7,3.748)(8,3.765)(9,3.773)(10,3.778)(11,3.780)(12,3.781)(13,3.782)
\uput[u](12.8,0){$n$}\uput[r](0,4.4){$u_n$}
\end{pspicture*}


\end{center}

\item Démontrer par récurrence que pour tout entier naturel $n$~:

\[2 \leqslant u_n \leqslant u_{n+1} \leqslant 4.\]

\item En déduire que la suite $(u_n)$ est convergente.
\end{enumerate}

\subsection*{Partie B}

On note $\ell$ la limite de la suite $(u_n)$.

Le but de cette partie est de trouver une valeur approchée de la limite $\ell$.

On considère la fonction $g$ définie et dérivable sur $[0~;~+\infty[$ par~:

\[g(x) = \ln(3x^2+1) - x.\]

\begin{enumerate}
\item On note $g'$ la fonction dérivée de la fonction $g$.

Démontrer que pour tout réel $x \in [0~;~+\infty[$,

\[g'(x) = \dfrac{-3x^2 + 6x - 1}{3x^2 + 1}.\]

\item Démontrer que la fonction $g$ est strictement décroissante sur $[2~;~4]$.
\item Démontrer que l'équation $g(x) = 0$ admet une unique solution $\alpha$ sur $[2~;~4]$ et en donner une valeur approchée à $10^{-2}$ près.
\item Justifier que $\ell = \alpha$.
\end{enumerate}

\bigskip

\end{document}
